\section{Formulario}

\subsection{Ley de Coulomb}

Para dos cargas puntuales \(q_1\) y \(q_2\) separadas una distancia \(r\), la magnitud de la fuerza eléctrica es

\[
\boxed{F=k\frac{|q_1q_2|}{r^2}}
\]

donde

\[
\boxed{k\approx\SI{8.99e9}{\newton\metre\squared\per\coulomb\squared}}
\]

y, para cálculos aproximados,

\[
\boxed{k\approx\SI{9.0e9}{\newton\metre\squared\per\coulomb\squared}}.
\]

La constante de Coulomb puede escribirse como

\[
\boxed{k=\frac{1}{4\pi\varepsilon_0}}
\]

con

\[
\boxed{\varepsilon_0\approx\SI{8.854e-12}{\farad\per\metre}}.
\]

Por tanto,

\[
\boxed{F=\frac{1}{4\pi\varepsilon_0}\frac{|q_1q_2|}{r^2}}.
\]

\begin{center}
\begin{tabular}{cl}
\toprule
Símbolo & Significado \\
\midrule
\(F\) & magnitud de la fuerza eléctrica \\
\(q_1,q_2\) & cargas eléctricas \\
\(r\) & distancia entre las cargas \\
\(k\) & constante de Coulomb \\
\(\varepsilon_0\) & permitividad eléctrica del vacío \\
\bottomrule
\end{tabular}
\end{center}

\begin{advertencia}
La expresión anterior se aplica directamente a cargas puntuales en reposo. También puede emplearse cuando los cuerpos pueden aproximarse razonablemente como cargas puntuales.

Con \(k\) expresada en unidades SI, las cargas deben estar en coulombs y las distancias en metros.
\end{advertencia}


\subsection{Carácter atractivo o repulsivo}

\[
\boxed{q_1q_2>0\Longrightarrow\text{repulsión}}
\]

\[
\boxed{q_1q_2<0\Longrightarrow\text{atracción}}.
\]

La fuerza actúa a lo largo de la recta que une las dos cargas.

\begin{advertencia}
En

\[
F=k\frac{|q_1q_2|}{r^2},
\]

\(F\) representa una magnitud y, por tanto, no es negativa. El signo de las cargas se utiliza para determinar el sentido de la interacción.
\end{advertencia}


\subsection{Dependencias de la fuerza}

\[
\boxed{F\propto |q_1q_2|}
\]

\[
\boxed{F\propto\frac{1}{r^2}}.
\]

Cambios frecuentes:

\[
q_1\to aq_1
\quad\Longrightarrow\quad
F\to aF
\]

si todo lo demás permanece constante.

\[
q_1\to aq_1,\qquad q_2\to bq_2
\quad\Longrightarrow\quad
F\to abF.
\]

\[
r\to ar
\quad\Longrightarrow\quad
F\to\frac{F}{a^2}.
\]

En particular,

\[
r\to2r\Longrightarrow F\to\frac{F}{4},
\]

\[
r\to3r\Longrightarrow F\to\frac{F}{9},
\]

\[
r\to\frac{r}{2}\Longrightarrow F\to4F.
\]


\subsection{Comparación entre dos situaciones}

Si las cargas o la distancia cambian,

\[
\boxed{
\frac{F'}{F}
=
\frac{|q_1'q_2'|}{|q_1q_2|}
\frac{r^2}{r'^2}
}.
\]

Esta relación permite comparar fuerzas sin calcularlas individualmente.


\subsection{Tercera ley de Newton}

Las fuerzas que dos cargas ejercen mutuamente satisfacen

\[
\boxed{\vec{F}_{1\to2}=-\vec{F}_{2\to1}}.
\]

Por tanto, sus magnitudes son iguales:

\[
\boxed{F_{1\to2}=F_{2\to1}}.
\]

Las fuerzas tienen sentidos opuestos, pero actúan sobre cuerpos diferentes.


\subsection{Forma vectorial de la ley de Coulomb}

Si \(\hat{\mathbf r}_{12}\) es un vector unitario dirigido desde \(q_1\) hacia \(q_2\),

\[
\boxed{
\vec{F}_{1\to2}
=
k\frac{q_1q_2}{r_{12}^2}\hat{\mathbf r}_{12}
}.
\]

En esta forma se conserva el signo de \(q_1q_2\).

Si

\[
q_1q_2>0,
\]

la fuerza tiene el mismo sentido que \(\hat{\mathbf r}_{12}\).

Si

\[
q_1q_2<0,
\]

la fuerza tiene sentido opuesto a \(\hat{\mathbf r}_{12}\).


\subsection{Forma vectorial mediante posiciones}

Si las cargas se encuentran en \(\vec r_1\) y \(\vec r_2\),

\[
\boxed{\vec r_{12}=\vec r_2-\vec r_1}
\]

\[
\boxed{r_{12}=|\vec r_2-\vec r_1|}
\]

\[
\boxed{
\hat{\mathbf r}_{12}
=
\frac{\vec r_2-\vec r_1}{|\vec r_2-\vec r_1|}
}.
\]

Entonces,

\[
\boxed{
\vec F_{1\to2}
=
kq_1q_2
\frac{\vec r_2-\vec r_1}
{|\vec r_2-\vec r_1|^3}
}.
\]


\subsection{Distancia entre cargas}

En una dimensión,

\[
\boxed{r=|x_2-x_1|}.
\]

En dos dimensiones,

\[
\boxed{
r=
\sqrt{(x_2-x_1)^2+(y_2-y_1)^2}
}.
\]

En tres dimensiones,

\[
\boxed{
r=
\sqrt{
(x_2-x_1)^2+
(y_2-y_1)^2+
(z_2-z_1)^2
}
}.
\]


\subsection{Principio de superposición}

La fuerza eléctrica neta sobre una carga \(q_i\) es

\[
\boxed{
\vec F_i=\sum_{j\neq i}\vec F_{j\to i}
}.
\]

Utilizando directamente las posiciones,

\[
\boxed{
\vec F_i
=
\sum_{j\neq i}
kq_iq_j
\frac{\vec r_i-\vec r_j}
{|\vec r_i-\vec r_j|^3}
}.
\]

Cada fuerza individual se calcula por separado y posteriormente se realiza la suma vectorial.


\subsection{Superposición en una dimensión}

Después de elegir un eje y un sentido positivo,

\[
\boxed{F_{\text{neta},x}=\sum_iF_{i,x}}.
\]

Por ejemplo, si \(+x\) apunta hacia la derecha,

\[
F_x>0\Longrightarrow\text{fuerza hacia la derecha},
\]

\[
F_x<0\Longrightarrow\text{fuerza hacia la izquierda}.
\]

\begin{advertencia}
El signo de una componente de fuerza depende de su sentido respecto del sistema de coordenadas, no del signo de la carga.
\end{advertencia}


\subsection{Componentes en dos dimensiones}

Si una fuerza de magnitud \(F\) forma un ángulo \(\theta\) medido desde el eje \(+x\),

\[
\boxed{F_x=F\cos\theta}
\]

\[
\boxed{F_y=F\sin\theta}.
\]

Los signos deben asignarse de acuerdo con el cuadrante del vector.


\subsection{Fuerza neta en dos dimensiones}

\[
\boxed{F_{\text{neta},x}=\sum_iF_{i,x}}
\]

\[
\boxed{F_{\text{neta},y}=\sum_iF_{i,y}}.
\]

En forma vectorial,

\[
\boxed{
\vec F_{\text{neta}}
=
F_{\text{neta},x}\hat{\mathbf i}
+
F_{\text{neta},y}\hat{\mathbf j}
}.
\]

Su magnitud es

\[
\boxed{
F_{\text{neta}}
=
\sqrt{
F_{\text{neta},x}^2+
F_{\text{neta},y}^2
}
}.
\]

Para determinar un ángulo de referencia,

\[
\boxed{
\tan\theta
=
\frac{|F_{\text{neta},y}|}
{|F_{\text{neta},x}|}
}.
\]

Después debe identificarse el cuadrante mediante los signos de las componentes.


\subsection{Caso de dos fuerzas con ángulo conocido}

Si dos fuerzas \(\vec F_1\) y \(\vec F_2\) forman entre sí un ángulo \(\theta\), la magnitud de la resultante puede calcularse mediante

\[
\boxed{
F_R=
\sqrt{
F_1^2+F_2^2+2F_1F_2\cos\theta
}
}.
\]

Esta expresión es equivalente a realizar la suma vectorial por componentes.

Casos particulares:

\[
\theta=0^\circ
\quad\Longrightarrow\quad
\boxed{F_R=F_1+F_2},
\]

\[
\theta=180^\circ
\quad\Longrightarrow\quad
\boxed{F_R=|F_1-F_2|},
\]

\[
\theta=90^\circ
\quad\Longrightarrow\quad
\boxed{F_R=\sqrt{F_1^2+F_2^2}}.
\]


\subsection{Equilibrio eléctrico}

Una carga se encuentra en equilibrio respecto de las fuerzas eléctricas cuando

\[
\boxed{\sum_i\vec F_i=\vec 0}.
\]

En componentes,

\[
\boxed{\sum_iF_{i,x}=0},
\qquad
\boxed{\sum_iF_{i,y}=0}.
\]

En una dimensión,

\[
\boxed{\sum_iF_{i,x}=0}.
\]

La fuerza neta nula no implica necesariamente ausencia de fuerzas individuales.


\subsection{Conversiones frecuentes}

\[
\boxed{\SI{1}{\milli\coulomb}=\SI{1e-3}{\coulomb}}
\]

\[
\boxed{\SI{1}{\micro\coulomb}=\SI{1e-6}{\coulomb}}
\]

\[
\boxed{\SI{1}{\nano\coulomb}=\SI{1e-9}{\coulomb}}
\]

\[
\boxed{\SI{1}{\pico\coulomb}=\SI{1e-12}{\coulomb}}.
\]


\subsection{Procedimiento rápido}

\begin{enumerate}
    \item Identificar la carga sobre la cual se pide la fuerza.
    \item Dibujar las cargas y establecer los ejes.
    \item Convertir cargas a coulombs y distancias a metros.
    \item Determinar si cada interacción es atractiva o repulsiva.
    \item Calcular cada magnitud con
    \[
    F=k\frac{|q_1q_2|}{r^2}.
    \]
    \item Asignar dirección y sentido a cada vector.
    \item Descomponer en componentes si es necesario.
    \item Aplicar
    \[
    \vec F_{\text{neta}}=\sum_i\vec F_i.
    \]
    \item Calcular magnitud y dirección de la resultante.
\end{enumerate}


\subsection{Resumen de fórmulas}

\begin{align*}
F &= k\frac{|q_1q_2|}{r^2}, \\[0.3em]
k &= \frac{1}{4\pi\varepsilon_0}, \\[0.3em]
\vec F_{1\to2} &= k\frac{q_1q_2}{r_{12}^2}\hat{\mathbf r}_{12}, \\[0.3em]
\vec F_{1\to2} &=
kq_1q_2\frac{\vec r_2-\vec r_1}{|\vec r_2-\vec r_1|^3}, \\[0.3em]
\vec F_i &= \sum_{j\neq i}\vec F_{j\to i}, \\[0.3em]
F_x &= F\cos\theta, \qquad F_y=F\sin\theta, \\[0.3em]
F_{\text{neta}} &=
\sqrt{F_{\text{neta},x}^2+F_{\text{neta},y}^2}, \\[0.3em]
\sum_i\vec F_i &= \vec 0 \qquad \text{(equilibrio)}.
\end{align*}